% !TEX root = ../main.tex
% !TEX program = lualatex
\documentclass[../main.tex]{subfiles}
\IfSubfilesClassLoaded{\externaldocument{\subfix{../build/main}}}{}
% =============================================================================
% File: 22_AIT_Fundamentals_and_Lessons_Learned.tex
% =============================================================================
\begin{document}
\IfSubfilesClassLoaded{\chapter{EDO Study Guide}}{}
\section{AIT Fundamentals and Lessons Learned}

\subsection{Summary}
\begin{description}
  \item[\textbf{Definition.}] An \ac{ait} is a trained team under an \ac{ait} manager responsible for installing and completing ship changes and alterations on specified ships~\autocite{edo-4-1-4-ait-fundamentals-2025}.
  \item[\textbf{Why use AITs.}] \acp{ait} apply specialized skills, leverage lessons learned across installations, and execute large numbers of ship changes efficiently~\autocite{edo-4-1-4-ait-fundamentals-2025}.
  \item[\textbf{Governance.}] \ac{ait} alteration development and execution follow the Navy Modernization Process Management and Operations Manual and TS9090-310 procedures aligned with the \ac{jfmm}~\autocite{edo-4-1-4-ait-fundamentals-2025}.
  \item[\textbf{Key roles.}] Sponsors, \ac{rmmco}, \ac{ait} managers, \ac{osic}, \ac{lma}, and \ac{nsa} coordinate planning, execution, and certification during availabilities~\autocite{edo-4-1-4-ait-fundamentals-2025}.
  \item[\textbf{MOA discipline.}] A \ac{moa} is required to align stakeholder expectations and support services before installation begins~\autocite{edo-4-1-4-ait-fundamentals-2025}.
\end{description}

\subsection{Alteration Installation Team}
\begin{description}
  \item[\textbf{Definition.}] An \ac{ait} is a unit (military, government activity, and/or contractor) trained and equipped to accomplish specific ship changes or alterations on specified ships~\autocite{edo-4-1-4-ait-fundamentals-2025}.
  \item[\textbf{Responsibility.}] \acp{ait} are responsible for the installation, performance, and completion of the assigned ship change or alteration~\autocite{edo-4-1-4-ait-fundamentals-2025}.
  \item[\textbf{Composition.}] Teams may be government-led, contractor-led, or a combination of government and contractor personnel~\autocite{edo-4-1-4-ait-fundamentals-2025}.
\end{description}

\subsection{Purpose and Advantages}
\begin{itemize}
  \item Use \acp{ait} when technical complexity or specialized oversight is required~\autocite{edo-4-1-4-ait-fundamentals-2025}.
  \item Reuse the same team to capture lessons learned and improve installation quality~\autocite{edo-4-1-4-ait-fundamentals-2025}.
  \item Apply \acp{ait} to large numbers of ships when repeatable alterations are required~\autocite{edo-4-1-4-ait-fundamentals-2025}.
  \item Example \ac{ait} alterations include \ac{canes}, \ac{aws}, and \ac{gpnts}~\autocite{edo-4-1-4-ait-fundamentals-2025}.
\end{itemize}

\subsection{Roles and Responsibilities}
\begin{description}
  \item[\textbf{AIT sponsor.}] Assigns tasks and funding to the \ac{ait} manager; sponsors may include \ac{syscom}, \ac{peo}, \ac{parm}, \ac{spm}, fleet commanders, \ac{tycom}, or \ac{cno} organizations~\autocite{edo-4-1-4-ait-fundamentals-2025}.
  \item[\textbf{AIT manager.}] Assigns and manages the \ac{ait}, designates the \ac{osic}, and plans and oversees alteration accomplishment in accordance with modernization policy and procedures~\autocite{edo-4-1-4-ait-fundamentals-2025}.
  \item[\textbf{AIT OSIC.}] Acts with the authority of the \ac{ait} manager, manages installation conduct, coordinates with the \ac{nsa}, and resolves issues and quality discrepancies~\autocite{edo-4-1-4-ait-fundamentals-2025}.
  \item[\textbf{AIT.}] Executes the alteration under the direction of the \ac{ait} manager or designated agent~\autocite{edo-4-1-4-ait-fundamentals-2025}.
  \item[\textbf{NSA.}] Integrates and verifies work accomplished during an availability, controls \ac{ait} access, ensures authorized work and compliance with technical specifications, and executes quality assurance monitoring~\autocite{edo-4-1-4-ait-fundamentals-2025}.
  \item[\textbf{LMA.}] Integrates all maintenance and modernization activities, coordinates testing controls, and reports work status to the \ac{nsa}~\autocite{edo-4-1-4-ait-fundamentals-2025}.
  \item[\textbf{RMMCO.}] Serves as the initial point of entry and gatekeeper for waterfront \ac{ait} installations, ensuring compliance with TS9090-310 requirements and delivery of logistics products for fleet self-sufficiency~\autocite{edo-4-1-4-ait-fundamentals-2025}.
\end{description}

\subsection{Need for RMMCO}
\begin{itemize}
  \item Rapid 1990s installation of command, control, communications, computers, intelligence, combat system, and hull, mechanical, and electrical improvements created coordination, interoperability, configuration control, and supportability challenges~\autocite{edo-4-1-4-ait-fundamentals-2025}.
  \item Many alterations lacked approved drawings and adequate integrated logistics support products~\autocite{edo-4-1-4-ait-fundamentals-2025}.
  \item Fleet leadership halted uncoordinated installations in 1998 and established \acp{rmmco} in 1999 to enforce integration~\autocite{edo-4-1-4-ait-fundamentals-2025}.
\end{itemize}

\subsection{Requirements}
\begin{description}
  \item[\textbf{AIT.}] The \ac{osic} coordinates with the \ac{nsa}, installs on a not-to-interfere basis, allows Ship's Force quality monitoring under the \ac{jfmm}, and follows an approved quality plan~\autocite{edo-4-1-4-ait-fundamentals-2025}.
  \item[\textbf{NSA.}] Certifies all work in the availability, integrates \ac{ait} work with other activity work, and coordinates with the \ac{lma} for quality monitoring~\autocite{edo-4-1-4-ait-fundamentals-2025}.
  \item[\textbf{Ship's Force.}] Maintains overall responsibility for the safety and security of \ac{ait} work~\autocite{edo-4-1-4-ait-fundamentals-2025}.
\end{description}

\subsection{Planning}
\begin{itemize}
  \item Generate the alteration list using ship program manager letters of authorization, \ac{tycom} authorizations, and inputs from \ac{parm}, \ac{navwar}, and \ac{navsea} 08 as required~\autocite{edo-4-1-4-ait-fundamentals-2025}.
  \item Review alteration scope using the \ac{scd} and validate impact on sequencing and availability integration~\autocite{edo-4-1-4-ait-fundamentals-2025}.
  \item Validate \ac{ait} quality systems; \ac{navsea} maintains the list of approved contractor and government quality systems~\autocite{edo-4-1-4-ait-fundamentals-2025}.
  \item The \ac{ait} manager coordinates early with the \ac{nsa} and provides production and test schedules before the availability starts~\autocite{edo-4-1-4-ait-fundamentals-2025}.
  \item \ac{ait} managers submit detailed support service requirements to the \ac{nsa} to integrate work, especially during private availabilities~\autocite{edo-4-1-4-ait-fundamentals-2025}.
  \item \ac{rmmco} check-in occurs 30 days prior to installation and requires authorization, integrated logistics support certification or risk assessment, test procedures, approved installation drawings, \acp{scd}, and \ac{ndenm} scheduling confirmation~\autocite{edo-4-1-4-ait-fundamentals-2025}.
  \item Conduct an \ac{ait} in-brief to reiterate requirements, deconflict support services, and confirm \ac{moa} expectations~\autocite{edo-4-1-4-ait-fundamentals-2025}.
\end{itemize}

\subsection{Memorandum of Agreement}
\begin{description}
  \item[\textbf{Purpose.}] \acp{ait} are not fully self-sufficient; the \ac{moa} aligns \ac{lma}, \ac{ait} managers, Ship's Force, and \ac{nsa} expectations and responsibilities~\autocite{edo-4-1-4-ait-fundamentals-2025}.
  \item[\textbf{Required content.}] Occupational safety, work control and tag-out, work hours, configuration management, quality assurance, cleanliness, support services, clearances, training, and fire safety protocols~\autocite{edo-4-1-4-ait-fundamentals-2025}.
  \item[\textbf{Timing.}] The \ac{moa} is signed out by A-1 to confirm stakeholder alignment before availability execution~\autocite{edo-4-1-4-ait-fundamentals-2025}.
\end{description}

\subsection{Execution}
\begin{description}
  \item[\textbf{Workmanship.}] The \ac{osic} enforces \ac{navsea} standard items and maintenance standards; planning yards provide oversight when tasked; housekeeping is required~\autocite{edo-4-1-4-ait-fundamentals-2025}.
  \item[\textbf{Testing.}] \ac{sovt} demonstrates operational capability and is included in integrated test plans, including \ac{subsafe} and security requirements~\autocite{edo-4-1-4-ait-fundamentals-2025}.
  \item[\textbf{Logistics turnover.}] Training, documentation, drawings, and allowance parts lists are turned over to the ship, along with spares and special test equipment~\autocite{edo-4-1-4-ait-fundamentals-2025}.
  \item[\textbf{Out-brief.}] The \ac{ait} conducts an out-brief with Ship's Force, \ac{nsa}, \ac{rmmco}, and \ac{tycom}; the completion report records status, deficiencies, and command comments~\autocite{edo-4-1-4-ait-fundamentals-2025}.
\end{description}

\subsection{Product Line Integration Walkthrough}
The path for a new equipment product line from production to integration into a ship or submarine involves the following standard steps:
\begin{enumerate}
  \item \textbf{System Production}: The manufacturer produces and packages the hardware and software components.
  \item \textbf{Factory Acceptance Testing (FAT)}: Verifies that the manufactured system meets design specifications and quality standards before shipment.
  \item \textbf{Staging and Delivery}: The system is shipped to a dedicated staging facility, warehouse, or directly to the \ac{isea} or installation site.
  \item \textbf{NMP and SCD Approval}: Under the Navy Modernization Process (NMP), a Ship Change Document (\ac{scd}) must be formally approved. The \ac{scd} details the physical change, weight and moment impacts, and power/cooling requirements, and is signed by the Program Manager and Technical Warrant Holder.
  \item \textbf{Installation Package Development}: The \ac{isea} or Planning Yard develops the Ship Installation Drawings (SIDs) and drafts the \ac{sovt} procedures.
  \item \textbf{RMMCO Gatekeeping}: The \ac{ait} Manager checks in with the Regional Maintenance Modernization Coordinating Office (\ac{rmmco}) (typically 30 days prior to installation) to verify the approved \ac{scd}, final SIDs, \ac{sovt} plan, and certified Integrated Logistics Support (ILS) products.
  \item \textbf{Waterfront Coordination}: The \ac{ait} coordinates with the \ac{nsa} and \ac{lma} to gain ship access. A Memorandum of Agreement (\ac{moa}) is signed by A-1 to align support services.
  \item \textbf{Physical Installation}: The \ac{ait} performs the physical installation on board (e.g., running cables, mounting racks, structural modifications).
  \item \textbf{System Testing}: The \ac{ait} and Ship's Force execute the \ac{sovt} (including cold and hot checks) to demonstrate operational capability.
  \item \textbf{Quality Assurance and Certification}: The \ac{nsa} and \ac{ait} \ac{osic} conduct final quality inspections and sign the certification/completion paperwork.
  \item \textbf{Logistics Handover}: The \ac{ait} turns over technical manuals, spares, training, and Allowance Parts Lists (APLs) to the ship's force.
  \item \textbf{Configuration Update}: The ship's configuration database (CDMD-OA) is updated within 30 days of completion to reflect the install.
\end{enumerate}

\subsection{Modernization Differences by Class}
Modernization processes differ significantly between ship classes depending on their nuclear or safety boundaries:
\begin{description}
  \item[\textbf{Conventional Surface Ships.}] Managed under the standard NMP. Changes are driven by \acp{scd} and gated on the waterfront by the local \ac{rmmco} to ensure logistics readiness before install.
  \item[\textbf{Submarines and Nuclear Vessels.}] Experience extremely tight configuration control due to \ac{subsafe}, Fly-By-Wire, and Nuclear Power (NAVSEA 08) boundaries:
    \begin{itemize}
      \item \textbf{Boundary Control}: Any change within a certified boundary (e.g., \ac{subsafe} or Fly-By-Wire) requires formal engineering authorization, execution by certified activities (e.g., public shipyards), and rigorous quality audits.
      \item \textbf{Nuclear Systems}: Alterations affecting the reactor plant or reactor support systems require direct, explicit NAVSEA 08 review and authorization.
      \item \textbf{Temporary Alterations (TempAlts)}: Submarine TempAlts are strictly time-bound, logged in active status lists, and must be completely uninstalled before the approved expiration date.
    \end{itemize}
  \item[\textbf{ISEA Involvement.}] The \ac{isea} serves as the lifecycle engineering agent for the system, responsible for developing the engineering changes, verifying SIDs, authoring \ac{sovt} plans, ensuring ILS product availability, and often providing the onsite \ac{ait} Manager or technical oversight.
\end{description}

\subsection{Lessons Learned}
\begin{itemize}
  \item Integration issues arise when \acp{ait}, Ship's Force, and industrial activities compete for the same space, systems, or services~\autocite{edo-4-1-4-ait-fundamentals-2025}.
  \item Validate the \ac{ait} quality system, ensure authorization is complete, and keep support service requirements current throughout execution~\autocite{edo-4-1-4-ait-fundamentals-2025}.
  \item Provide \ac{ait} production schedules early, sign the \ac{moa} before the availability, and conduct in-brief and out-brief sessions~\autocite{edo-4-1-4-ait-fundamentals-2025}.
\end{itemize}

\subsection{Quick Review}
\begin{itemize}
  \item \acp{ait} are used when specialized work warrants, installations must be repeated, or many ships require the same change~\autocite{edo-4-1-4-ait-fundamentals-2025}.
  \item \ac{ait} installations may be performed by government teams, contractors, or a combination~\autocite{edo-4-1-4-ait-fundamentals-2025}.
  \item \ac{rmmco} serves as the gatekeeper for waterfront alteration installations~\autocite{edo-4-1-4-ait-fundamentals-2025}.
  \item The \ac{moa} promotes stakeholder discussion and alignment during planning~\autocite{edo-4-1-4-ait-fundamentals-2025}.
\end{itemize}

\IfSubfilesClassLoaded{
  \RenewDocumentCommand{\entryname}{}{\textbf{\color{Modern} Acronym}}
  \RenewDocumentCommand{\descriptionname}{}{\textbf{\color{Modern} Definition}}
  \printnoidxglossary[
    type=\acronymtype,
    title=Acronyms,
    style=long-booktabs
  ]}{}
\end{document}
